% ============================================================================
%  Chapter 2: Related Work
% ============================================================================
\section{Related Work}
\label{sec:related-work}

\subsection{Decentralized identity and verifiable credentials}
The W3C has standardized Decentralized Identifiers (DIDs) \cite{w3c-did} and Verifiable
Credentials (VCs) \cite{w3c-vc} as mutually independent, composable primitives: a DID is a
URI that resolves to a DID document carrying control and verification methods, and a VC is a
cryptographically verifiable data model for claims such as qualifications, authorization, and
reputation. Implementation layers such as ION \cite{ion} anchor DID operations on Bitcoin without
a new token. These standards provide the skeleton for issuing identities to agents, but they do not
by themselves determine how identities are trusted in commerce, how disputes are adjudicated, or
how reputation is priced --- the concerns of this paper.

\subsection{Soulbound tokens and bound NFTs}
Soulbound tokens (EIP-5192 \cite{eip-5192}) and account-bound NFTs (EIP-4973 \cite{eip-4973})
generalize non-transferability, letting an address hold verifiable, non-tradeable attestations of
membership, achievement, or commitment. In our architecture the agent identity (an ERC-721 in
\textsf{AgentRegistry}) is bound to a real principal at registration and cannot be detached from the
responsibility it carries --- the accountability analogue of the soulbound idea.

\subsection{Autonomous accounts and payments}
Account abstraction (ERC-4337 \cite{eip-4337}) separates the signing key from the account logic,
enabling agents to sponsor gas, batch operations, and authorize delegated execution. The
agent-to-agent finance literature \cite{a2a-finance} surveys blockchain payment and trust
infrastructure for autonomous agents. Our work assumes such an account layer as the substrate and
focuses on the higher-level question of how two accounts that have never met reach trust: guarantee,
adjudicate, and build reputation.

\subsection{On-chain reputation and attestation}
The Ethereum Attestation Service (EAS) \cite{eas} provides a general on-chain attestation
registry. AgentReputation \cite{agent-rep} proposes a decentralized reputation framework for
agentic AI. Recent datasets of blockchain-registered AI agents \cite{agent-dataset} indicate early
adoption of agent identity standards. These works establish that reputation \emph{can} be recorded
on-chain, but leave open who is authorized to write, how truthful attestations are incentivized, and
how reputation is priced into risk. Our reputation hub (Section~\ref{sec:contracts}) closes the
write-incentive gap by restricting writers to audited escrow/voting contracts (no self-rating) and
by coupling scores to guarantee pricing.

\subsection{Decentralized arbitration and guarantees}
Decentralized dispute resolution has been explored by Kleros \cite{kleros} (a jury selected from a
set of paid, random jurors), Aragon Court \cite{aragon} (aragon.org/dish), and UMA's optimistic
oracle \cite{uma} (price oracle with a challenge period). Kleros and Aragon Court select jurors by
\emph{lottery} and reward coherent majorities; UMA incentivizes truthful price reporting through
disputes and forfeits. These mechanisms share the Schelling-point insight --- that truthful
majorities are self-enforcing --- but none of them \emph{closes the loop} with reputation and
guarantee pricing: a Kleros verdict is final and does not update a merchant's future cost of doing
business. On the insurance side, Nexus Mutual \cite{nexus-mutual} runs a mutual cover pool with a
claims-assessment community. Our contribution is to make the guarantee itself (escrow plus
guarantor stake) the stake that voters price, and to make every verdict feed a reputation score that
prices the next guarantee --- a three-way closed loop (Section~\ref{sec:pricing}).

\subsection{Trust research specific to agents}
Recent work on trustworthy LLM agents \cite{agent-networks, zero-trust, insured-agents} and on
binding agent identities \cite{binding-id} converge on the need for attestation, insurance, and
verifiable accountability. ERC-8004 (trustless agents) \cite{erc8004} sketches an interoperability
standard. We position this paper as the first to combine, end-to-end and with a formal
incentive analysis, (i) a registered identity, (ii) a guaranteed escrow, (iii) staked-voting
arbitration, and (iv) reputation-driven guarantee pricing in a single protocol whose contracts and
empirical validation are public.
